% Part V -- AMD case study and quantitative analysis (target: about 25 pages)
% Method: Option A (analysis of published research and benchmarks); no new
% hardware measurements and no simulator.
\clearpage
\part{AMD Case Study and Quantitative Analysis}
\label{part:case-study}
\partbudget{25}{method, AMD-reported evidence and confound grid drafted; independent evidence, Zen~2 to Zen~3 comparison and summary added}

\section{From Zen 2 Chiplets to Zen 3 and 3D V-Cache}
\label{sec:zen-progression}

\input{"../initial-draft-2/figures/zen-ccd-evolution.tex"}

\Cref{tab:zen-progression} and \cref{fig:zen-ccd-evolution} summarise the changes in cache organisation that matter for this paper, using only AMD's published descriptions.

\begin{table}[htbp]
  \centering
  \caption{L3 organisation per compute die across the generations studied (AMD-published figures)}
  \label{tab:zen-progression}
  \small
  \begin{tabularx}{\textwidth}{@{}>{\raggedright\arraybackslash}p{3.3cm}>{\raggedright\arraybackslash}X>{\raggedright\arraybackslash}X>{\raggedright\arraybackslash}X@{}}
    \toprule
    & \textbf{Zen 2 CCD} & \textbf{Zen 3 CCD} & \textbf{Zen 3 CCD with 3D V-Cache} \\
    \midrule
    Cores per CCD & 8 & 8 & 8 \\
    CCX organisation & Two 4-core CCXs & One 8-core CCX & One 8-core CCX \\
    L3 shared by one core & Up to 16 MB (its CCX) & 32 MB & 96 MB \\
    L3 per CCD & 2 $\times$ 16 MB & 32 MB & 32 MB base + 64 MB stacked \\
    Private L2 per core & 512 KB & 512 KB & 512 KB \\
    CCD area & 74 mm$^2$ & 81 mm$^2$ & 81 mm$^2$ + 41 mm$^2$ L3D \\
    Added L3 latency & --- & --- & +4 cycles \\
    Sources & \cite{suggs2019zen2,naffziger2020chiplet} & \cite{amd2022epyc7003,wuu2022vcache} & \cite{amd2022epyc7003,wuu2022vcache} \\
    \bottomrule
  \end{tabularx}
\end{table}

Two steps stand out. Zen~3 did not add L3 capacity per CCD but changed its \emph{sharing}: each core can now use the whole 32\,MB rather than only its own CCX's 16\,MB \cite{amd2022epyc7003,suggs2019zen2}. 3D V-Cache then tripled the capacity seen by each core without changing the CCD's footprint or the sharing structure \cite{amd2022epyc7003,wuu2022vcache}. These two steps give a natural experiment for the effect of L3 capacity, as long as other differences are controlled (\cref{sec:confounds}).

\subsection{Zen~2 to Zen~3: a same-core-count comparison}
\label{sec:zen2-zen3}
The Zen~2 to Zen~3 step can be examined with two eight-core, single-CCD desktop parts. The Ryzen~7 3700X (Zen~2) and Ryzen~7 5800X (Zen~3) both have eight cores, sixteen threads and 32\,MB of L3, but the 3700X splits its L3 into two 16\,MB halves, one per four-core CCX, while the 5800X shares all 32\,MB among its eight cores \cite{suggs2019zen2,amd2022epyc7003}. They differ in other ways too: base and maximum boost clocks of 3.6 and 4.4\,GHz against 3.8 and 4.7\,GHz, and a TDP of 65\,W against 105\,W \cite{advancedmicrodevicesinc.AMDRyzen30002022}.

AMD attributes a 19\% IPC increase to Zen~3, citing the unified eight-core complex with direct access to 32\,MB of L3 among the core improvements. The figure comes from 25 workloads run at a locked 4\,GHz on the eight-core Ryzen~7 3800XT and 5800X, both with DDR4-3600 memory \cite[note~4]{advancedmicrodevicesinc.AMDLaunchesAMD2020}. Because clock, core count and memory are held equal, this is a cleaner comparison than most vendor figures. It still cannot separate the shared L3 from the many other changes to the Zen~3 core.

Independent testing of the two retail parts shows the combined effect. With the same DDR4-3200 CL14 memory, TechSpot measured the 5800X ahead of the 3700X by 22\% in Cinebench R20 multi-core, 24\% and 19\% in 7-Zip compression and decompression, 16\% in Blender, 30\% in V-Ray, 26\% in Corona, 12\% in code compilation, 6\% in DaVinci Resolve, 14\% in Premiere Pro, 27\% in Photoshop and 23\% in After Effects \cite{waltonAMDRyzen72020}. During the multi-core Cinebench run the 5800X held about 4.55\,GHz, and it drew about 20\,W more than the 3700X \cite{waltonAMDRyzen72020}. These are whole-processor differences. The 5800X runs at higher clocks under a higher power limit and has a redesigned core, so the gains cannot be credited to the unified L3 alone. The gains are broad rather than concentrated in cache-sensitive work: the largest is in V-Ray, a multi-threaded renderer (30\%), and large gains also appear in Photoshop (27\%), which the reviewer describes as relying heavily on single-core performance \cite{waltonAMDRyzen72020}. In our reading, this points to the core and clock as much as to the cache. Cache measurements give a more direct view. The FPS Review ran AIDA64's cache and memory benchmark on both processors with identical DDR4-3600 CL16 memory and the same fabric clock (1800\,MHz) \cite{justiceAMDRyzen72021}. Memory read bandwidth was the same, about 50.7\,GB/s on both, and memory latency fell from 68.8\,ns on the 3700X to 61.8\,ns on the 5800X. L3 read bandwidth was also nearly identical (633 against 631\,GB/s). L3 latency, however, \emph{rose} from 9.9\,ns to 10.9\,ns \cite{justiceAMDRyzen72021}. At the core clocks the tool reported during the runs (4.40 and 4.85\,GHz), this is roughly 44 cycles against 53 cycles. The figures come from screenshots of a single run on each processor, so the cycle counts are only approximate. Chips and Cheese's latency measurements across Zen generations agree: unifying the two 16\,MB caches of a Zen~2 CCD into one 32\,MB cache added 7 to 8 cycles of L3 latency, whereas doubling the size of each L3 slice from Zen to Zen~2 added about 5 cycles and tripling it with V-Cache only 3 to 4, which suggests that enlarging slices is cheaper in latency than merging separate caches into one \cite{cozmaDeepDivingZen2022}. \Cref{fig:zen2-zen3-latency} plots the AIDA64 latencies.

\input{"../initial-draft-2/figures/zen2-zen3-latency.tex"}


This result fits the trade-off described in \cref{sec:planar-limits}. Doubling the cache a core can reach, from 16 to 32\,MB, removes misses for working sets between the two sizes and the traffic between CCXs, but a larger shared array serving more cores takes longer to reach. The 5800X's faster single-threaded applications, in which the 32\,MB working-set limit rarely matters, therefore most likely owe their gains to the core and clock, not to faster cache access. Isolating the benefit of the unified L3 itself would require fixed-clock measurements of workloads with working sets between 16 and 32\,MB, or of communication between cores in different CCXs, which none of the sources in hand provides.

Two later generations extend the same pattern (\cref{tab:later-x3d}). With Zen~4, AMD kept the stack unchanged in principle but moved the CCD to 5\,nm while the cache die stayed at 7\,nm \cite{alam2ndGeneration3D2024}; on servers, up to twelve CCDs per package each reach 96\,MB of L3 with V-Cache, for up to 1152\,MB per socket \cite{advancedmicrodevicesinc.4thGenAMD2023}. With Zen~5, AMD inverted the stack so that the cache die sits below the CCD (\cref{sec:thermal-evidence}).

\begin{table}[htbp]
  \centering
  \caption{3D V-Cache across Zen generations (AMD-published figures unless noted)}
  \label{tab:later-x3d}
  \small
  \begin{tabularx}{\textwidth}{@{}>{\raggedright\arraybackslash}p{2.0cm}>{\raggedright\arraybackslash}X>{\raggedright\arraybackslash}X>{\raggedright\arraybackslash}X@{}}
    \toprule
    & \textbf{Zen 3} & \textbf{Zen 4} & \textbf{Zen 5} \\
    \midrule
    Example parts & Ryzen~7 5800X3D; EPYC 7773X & Ryzen~7 7800X3D; EPYC 9684X & Ryzen~7 9800X3D \\
    CCD / cache die & 7\,nm / 7\,nm & 5\,nm / 7\,nm & 4\,nm (TSMC) / not stated \\
    Stack order & cache die on top of CCD & cache die on top of CCD & CCD on top of cache die \\
    L3 per V-Cache CCD & 96\,MB & 96\,MB & 96\,MB (32 + 64) \\
    Largest socket L3 & 768\,MB & 1152\,MB & --- \\
    Added L3 latency & 3--4 cycles (measured) & about 4 cycles (measured) & 4 cycles (measured) \\
    Sources & \cite{wuu2022vcache,amd2022epyc7003,cozmaDeepDivingZen2022} & \cite{alam2ndGeneration3D2024,advancedmicrodevicesinc.4thGenAMD2023,lamAMDs7950X3DZen2023} & \cite{advancedmicrodevicesinc.GamingLegendContinues2024,advancedmicrodevicesinc.AMDRyzen72024,lambertHeterogenousIntegrationHigh2025,cozmaAMDs9800X3D2nd2024} \\
    \bottomrule
  \end{tabularx}
\end{table}
AMD lists the Ryzen~7 9800X3D with 96\,MB of L3, cores made in TSMC's 4\,nm process and an I/O die in 6\,nm \cite{advancedmicrodevicesinc.AMDRyzen72024}; it does not state the process used for the cache die.

\section{Comparing Architectural Configurations}
\label{sec:config-compare}
The comparisons in this part rely on pairs of processors that differ mainly in L3 capacity. \Cref{tab:config-pairs} lists the pairs used and what else differs between them.

\begin{table}[htbp]
  \centering
  \caption{Processor pairs that differ mainly in L3 capacity}
  \label{tab:config-pairs}
  \small
  \begin{tabularx}{\textwidth}{@{}>{\raggedright\arraybackslash}p{3.3cm}>{\raggedright\arraybackslash}p{1.1cm}>{\raggedright\arraybackslash}p{2.2cm}>{\raggedright\arraybackslash}X>{\raggedright\arraybackslash}p{1.6cm}@{}}
    \toprule
    \textbf{Pair} & \textbf{Cores} & \textbf{L3} & \textbf{Other differences} & \textbf{Source} \\
    \midrule
    EPYC 7763 vs 7773X (Zen~3) & 64 & 256 vs 768\,MB & base clock 2.45 vs 2.20\,GHz; same 3.5\,GHz boost & \cite{larabelAMDEPYC7773X2022} \\
    Ryzen~7 5800X vs 5800X3D (Zen~3) & 8 & 32 vs 96\,MB & 5800X clocks slightly higher & \cite{larabelAMDRyzen72022} \\
    7950X3D: plain CCD vs V-Cache CCD (Zen~4) & 8 per CCD & 32 vs 96\,MB & plain CCD about 7\% faster at stock; identical when both are fixed at 4.2\,GHz & \cite{lamAMDs7950X3DZen2023} \\
    EPYC 9654 vs 9684X (Zen~4) & 96 & 384 vs 1152\,MB & all-core boost 3.55 vs 3.42\,GHz; base 2.40 vs 2.55\,GHz; same 3.70\,GHz maximum; default power 360 vs 400\,W & \cite{advancedmicrodevicesinc.AMDEPYC90042023} \\
    \bottomrule
  \end{tabularx}
\end{table}

\section{Method for Analysing Published Evidence}
\label{sec:evidence-method}

The quantitative analysis uses published research and benchmarks only (\cref{fig:evidence-method}). This is the most manageable option given that the group does not have access to the relevant hardware. Its limitation is dependence on the experimental conditions and measurements that others chose to publish; the analysis can only be as controlled as its sources.

\input{"../initial-draft-2/figures/evidence-method.tex"}

Each source is recorded with: the processors compared; clock frequencies (fixed or boost); memory type, speed and timings; power limits; software and settings; the metric (for example, frames per second, IPC, runtime or throughput); and whether the source is the vendor or an independent party. Results are reported as relative changes with the source, metric and conditions stated, never as a single headline number.

\textbf{Source selection.} Sources are admitted in three tiers: peer-reviewed papers; independent benchmark publications that state their hardware, settings and method; and vendor material, which is used only when labelled as such. Every source used must have its PDF in the project's source folder. Where sources conflict, the one with the better-controlled comparison (\cref{sec:confounds}) is preferred and the disagreement is reported. Averages across workloads are formed as geometric means of performance ratios, so that no single workload dominates.

\section{Published Benchmarks and Performance Measurements}
\label{sec:benchmarks}

The evidence falls into two groups: AMD-reported results (\cref{tab:amd-reported}) and independent measurements (\cref{tab:independent}).

\begin{table}[htbp]
  \centering
  \caption{AMD-reported performance figures relevant to this study (vendor claims, not independent measurements)}
  \label{tab:amd-reported}
  \small
  \begin{tabularx}{\textwidth}{@{}>{\raggedright\arraybackslash}X>{\raggedright\arraybackslash}p{4.4cm}>{\raggedright\arraybackslash}p{2.0cm}@{}}
    \toprule
    \textbf{Claim} & \textbf{Conditions stated} & \textbf{Source} \\
    \midrule
    About 15\% faster gaming, 8-core Ryzen~7 5800X3D vs 12-core Ryzen~9 5900X; per-title results from a tie to 1.36$\times$ & Six games at 1080p, high preset; both with DDR4-3600 and RTX~3080; AMD testing, Dec.\ 2021 & \cite[slides~29, 35]{wuu2022vcache} \\
    About 66\% more RTL-simulation jobs per hour (24.4 to 40.6), 16-core EPYC with V-Cache vs 16-core EPYC~73F3 & Synopsys VCS; same reference platform; AMD testing, Sept.\ 2021 & \cite[slides~31, 35]{wuu2022vcache} \\
    15\% average gaming uplift, 12-core 3D-chiplet prototype vs Ryzen~9 5900X & Both at the same fixed frequency; disclosed games only & \cite{swaminathanAdvancedPackagingEnabling2021} \\
    15\% IPC improvement from Zen to Zen~2 & Many simultaneous changes, including doubled L3 & \cite{suggs2019zen2} \\
    19\% IPC improvement from Zen~2 to Zen~3 (one AMD test, reported in both sources) & 25 workloads at a fixed 4\,GHz, 3800XT vs 5800X, equal memory; but a redesigned core as well as the unified L3 & \cite{advancedmicrodevicesinc.AMDLaunchesAMD2020}, \cite[slides~10, 35]{wuu2022vcache} \\
    19\% average memory-latency reduction, Rome vs Naples & Central IOD, decoupled clocks, optimised IFOP & \cite{naffziger2020chiplet} \\
    \bottomrule
  \end{tabularx}
\end{table}

\begin{table}[htbp]
  \centering
  \caption{Independent measurements of the effect of 3D V-Cache}
  \label{tab:independent}
  \small
  \begin{tabularx}{\textwidth}{@{}>{\raggedright\arraybackslash}X>{\raggedright\arraybackslash}p{4.6cm}>{\raggedright\arraybackslash}p{1.5cm}@{}}
    \toprule
    \textbf{Result} & \textbf{Conditions} & \textbf{Source} \\
    \midrule
    IPC $+4.9\%$ to $+19\%$ across five workloads; one workload slightly negative & 7950X3D, each CCD at a fixed 4.2\,GHz; performance counters & \cite{lamAMDs7950X3DZen2023} \\
    Up to 3.4$\times$ on MiniFE for a subset of problem sizes & EPYC 7773X vs 7763, 64 cores each & \cite{domkeLocusPerformanceQuantifying2023} \\
    OpenFOAM from 0.81$\times$ to 1.54$\times$ (motorBike) and up to 1.47$\times$ (urban air pollution); L3 miss rate down by up to 30 points & EPYC 7773X vs 7763, dual socket, clocks not fixed & \cite{lawendaEvaluatingImpactL32025} \\
    gem5 $+7.6\%$, libx264 and 7-Zip about $+5\%$ on one CCD; Milan-X clocked about 5\% lower & Milan-X vs Milan cloud instances, clocks estimated & \cite{cozmaDeepDivingZen2022} \\
    Of 350 significant differences among 420 tests, 106 favour the 5800X3D (12 by 50\% or more) and 244 the 5800X (all but 3 by under 10\%) & Ryzen~7 5800X3D vs 5800X, Linux & \cite{larabelAMDRyzen72022} \\
    L3 latency $+3$ to $+4$ cycles & Milan-X, 7950X3D, 9800X3D & \cite{cozmaDeepDivingZen2022,lamAMDs7950X3DZen2023,cozmaAMDs9800X3D2nd2024} \\
    \bottomrule
  \end{tabularx}
\end{table}

The independent results agree with AMD's on direction but not on size: they show a wider spread, including small losses, and the best-controlled of them, Lam's fixed-frequency per-CCD measurements, shows single-digit to high-teen IPC gains rather than uniform double-digit speedups. \Cref{fig:vendor-vs-independent} shows the difference in spread.

\input{"../initial-draft-2/figures/vendor-vs-independent.tex"}


\subsection{Latency through the whole hierarchy}
\label{sec:milanx-latency}
The clearest picture of what the stacked cache changes comes from a latency sweep, in which a test program walks through a region of memory of increasing size and measures the time per load. Chips and Cheese ran such a sweep on Milan and Milan-X cloud instances (\cref{fig:milanx-latency}) \cite{cozmaAMDsVCacheTested2022}. The two curves coincide through L1 (4 cycles) and L2 (about 12 cycles). In the L3 range Milan-X is slightly slower, 51.14 against 47.49 cycles, the three to four extra cycles discussed in \cref{sec:vcache-access}. The curves separate at 32\,MB: Milan's latency climbs towards memory, while Milan-X's stays flat until beyond about 64\,MB and then rises as its 96\,MB L3 is exceeded. Between those sizes a load costs about 51 cycles on Milan-X, while on Milan it rises to roughly 250--300 cycles by 64--96\,MB, five to six times as long. Beyond 96\,MB both go to memory, at 342 cycles for Milan and 313 for Milan-X; in nanoseconds Milan-X's memory latency was also lower (89 against 106\,ns), and the article notes that Milan-X was boosting better than Milan in this test \cite{cozmaAMDsVCacheTested2022}.

A third source agrees. On a Zen~3 Ryzen~5 5600G, which has a 16\,MB L3 on a single die, 7-cpu.com measured 4 cycles for L1, 12 for L2, 47 for L3 and 47 cycles plus 74\,ns for memory, about 375 cycles at the chip's 4.45\,GHz \cite{7-cpu.comAMDZen3}. The same source shows a second effect that latency sweeps can mix up with cache misses: with ordinary 4\,KB pages, latency at 512\,KB was 21 cycles, against 12 with 2\,MB pages, although the data still fits in L2; the source's table attributes part of the rise to misses in the 64-entry first-level data TLB, which can no longer hold the address translations, and leaves the exact 4\,KB penalty unstated \cite{7-cpu.comAMDZen3}. Sweeps such as Chips and Cheese's therefore measure cache and TLB effects together, and the steps in \cref{fig:milanx-latency} are not purely cache capacities.

These measurements also check the textbook's numbers. The measured memory latency of 313--342 cycles lies just above the 100--300-cycle range assumed in \cref{sec:tradeoff-example}, so on these processors a removed miss saves at least as much as that example allowed. With the measured values, \cref{eq:vcache-tradeoff} gives a break-even of $3.65/313 \approx 1.2$ percentage points of L3 misses removed for Milan-X to come out ahead.

\input{"../initial-draft-2/figures/milanx-latency.tex"}

\subsection{A working-set sweep: MiniFE}
\label{sec:minife}
Domke \emph{et al.} ran the MiniFE finite-element proxy application on dual-socket Milan and Milan-X systems with the same 64 cores per socket, 280\,W TDP and 16-channel DDR4 memory, while scaling the problem from $100^3$ to $400^3$ grid points \cite{domkeLocusPerformanceQuantifying2023}. \Cref{fig:minife-sweep} shows the result: a sweep through working-set sizes of exactly the kind sketched in \cref{fig:working-set}.
\begin{itemize}
  \item At the smallest sizes ($n = 100$ and 110) Milan-X is about 10\% \emph{slower}. Both systems hold the data in L3, so the difference is probably Milan's higher base clock (2.45 against 2.20\,GHz); \cref{sec:clock-correction} tests this.
  \item From $n = 120$ the advantage rises steeply, to 3.4$\times$ at $n = 160$. In this range Milan's 512\,MiB of L3 (two sockets) no longer holds the data, while Milan-X's 1536\,MiB still does.
  \item Beyond $n = 170$ the advantage falls quickly and settles at about 1.04$\times$ from $n \approx 330$ onwards, once the data exceeds both caches.
\end{itemize}

\textbf{Worked example: inferring the working set.} The curve can be used to estimate how much data each grid point needs, with no information about MiniFE's data structures. Milan's own performance falls sharply between $n = 120$ and 130 \cite[Fig.~1]{domkeLocusPerformanceQuantifying2023}, that is, between $1.73$ and $2.20$ million points; if this is where the data outgrows 512\,MiB, each point needs about $512 \times 2^{20} / (1.73\text{ to }2.20 \times 10^6) \approx 245$--310 bytes. Milan-X's performance falls between $n = 170$ and 180, between $4.91$ and $5.83$ million points; with 1536\,MiB of L3 that gives about 275--330 bytes per point. The two independent estimates agree, which supports the reading that the speedup curve is a cache-capacity effect: about 300 bytes per grid point fit in three times as many points' worth of L3 on Milan-X. The estimate is approximate, because not all of the L3 holds useful data and the two sockets' caches are not perfectly shared.

\input{"../initial-draft-2/figures/minife-sweep.tex"}

\subsection{OpenFOAM in detail}
\label{sec:openfoam}
Lawenda \emph{et al.}'s OpenFOAM study adds measured miss rates and clock speeds \cite{lawendaEvaluatingImpactL32025}. For the dual-socket 64-core Rome, Milan and Genoa systems, all with 256\,MB of L3 per socket, the L3 miss rate rose from about 10--15\% on the smallest motorBike meshes to about 80\% on the largest. Milan-X, with 768\,MB per socket, cut the miss rate by up to 30 percentage points on the mid-sized to large meshes. Its performance relative to Milan ranged from 0.81$\times$ to 1.54$\times$ across the eight mesh sizes: slower on some small meshes, fastest at the ``mhigh'' mesh of 6.7 million cells, and falling back to 1.24$\times$ on the largest, 39.4-million-cell mesh, where the data exceeds even Milan-X's cache. Expressed per core, Lawenda \emph{et al.} place Milan-X's advantage between about 15\,000 and 100\,000 cells per core, although their Fig.~4 still shows it 1.22$\times$ ahead at about 308\,000 cells per core (xhigh); and its best utilisation came at about 14\,800 cells per core against about 4700 for Rome and Milan. The newer generation shows the same pattern: Genoa-X ranged from 0.98$\times$ to 1.60$\times$ against the 96-core Genoa \cite{lawendaEvaluatingImpactL32025}. \Cref{fig:openfoam-sweep} plots the ratios Lawenda \emph{et al.}\ label in their figure for the five larger meshes.

\input{"../initial-draft-2/figures/openfoam-sweep.tex"}


The study's measured clocks show the confound discussed in \cref{sec:confounds}: during the runs, Milan ran at 2.9--3.2\,GHz and Milan-X at 2.8--3.0\,GHz \cite[Table~5]{lawendaEvaluatingImpactL32025}. Milan-X's losses on the small meshes, where both systems' caches suffice, are therefore at least partly a clock effect, as in MiniFE.

\subsection{Breadth: 350 desktop results}
\label{sec:phoronix}
Phoronix compared the eight-core Ryzen~7 5800X3D and 5800X across 420 benchmarks on Linux and charted the 350 with a statistically significant difference \cite{larabelAMDRyzen72022}. Classifying each result by the colour of its bar in that chart gives the distribution in \cref{fig:phoronix-hist}. The 5800X was faster in 244 tests, but all except three of those margins were under 10\% and the median was about 4\%, consistent with its higher clocks. The 5800X3D was faster in 106 tests, with a median margin of about 11\%; 58 of its wins were 10\% or more and 12 exceeded 50\%, led by ASKAP's tConvolve gridding kernel (294\%), Zstd compression, oneDNN, CloverLeaf and OpenFOAM (68\%). Phoronix's own summary describes the 5800X winning ``a small set of tests''; the chart shows that the 5800X won more tests but by much smaller margins. The distribution is the clearest available picture of the workload dependence: for most desktop software the stacked cache costs a few percent, and for a minority it is worth tens of percent or more.

\input{"../initial-draft-2/figures/phoronix-hist.tex"}

\section{Cache-Sensitive Versus Cache-Insensitive Workloads}
\label{sec:workload-sensitivity}
The evidence sorts workloads into three groups, matching the working-set argument of \cref{sec:vcache-workloads}:
\begin{itemize}
  \item \textbf{Strongly cache-sensitive:} memory-bound HPC and engineering codes whose working sets fall between the base and stacked capacities (MiniFE, OpenFOAM at mid-sized meshes, RTL simulation) and games with high L3 miss rates. Gains range from tens of percent to more than 3$\times$ \cite{domkeLocusPerformanceQuantifying2023,lawendaEvaluatingImpactL32025,lamAMDs7950X3DZen2023}.
  \item \textbf{Mildly sensitive:} workloads with moderate miss rates, such as compression and video encoding, gain around 5--10\% at equal clock \cite{lamAMDs7950X3DZen2023,cozmaDeepDivingZen2022}, which a V-Cache part's lower clock can cancel.
  \item \textbf{Insensitive or negative:} workloads whose working sets already fit (small OpenFOAM meshes, the flight simulator in Lam's tests) or that are compute-bound gain nothing and can lose slightly from the added latency or lower clocks \cite{lawendaEvaluatingImpactL32025,lamAMDs7950X3DZen2023,larabelAMDRyzen72022}.
\end{itemize}
The useful predictor is therefore not the application type but its L3 miss rate on the base configuration: in Lam's results, the workloads with the most L3 misses gained most and those with the fewest gained least or lost \cite{lamAMDs7950X3DZen2023}.

\section{Separating Cache Effects From Other Factors}
\label{sec:confounds}

A difference in measured performance between two processors can rarely be attributed to cache size alone. Using \cref{eq:perf} and \cref{eq:amat}, the main confounding factors are:
\begin{itemize}
  \item \textbf{Clock frequency.} If the larger-cache part runs at a different frequency, $f$ in \cref{eq:perf} changes as well as IPC. Fixed-frequency comparisons, like the one AMD reports for its prototype \cite{swaminathanAdvancedPackagingEnabling2021}, remove this factor.
  \item \textbf{Core count and topology.} More cores or a different number of CCDs changes throughput and the cache available per core. AMD's headline desktop comparison sets an eight-core, single-CCD Ryzen~7 5800X3D against a twelve-core, two-CCD Ryzen~9 5900X \cite[slide~29]{wuu2022vcache}, so it mixes cache capacity with core count, CCD topology and clock speed. Its server comparison holds the core count at 16 but uses the EPYC~73F3 as the baseline, an F-series model that AMD positions for high frequency \cite{amd2022epyc7003,wuu2022vcache}.
  \item \textbf{Memory bandwidth and latency.} Memory speed, timings and the fabric clock change the miss penalty in \cref{eq:amat} and can make a larger cache appear more or less valuable.
  \item \textbf{Concurrent microarchitectural changes.} Generational comparisons mix many changes. AMD's 15\% Zen-to-Zen~2 IPC figure, for example, combines a new branch predictor, a larger op cache, wider floating-point paths and a doubled L3, among others \cite{suggs2019zen2}; it cannot be credited to the cache.
  \item \textbf{Power limits and software.} Package power limits, operating-system scheduling across CCDs, and application versions all affect results.
\end{itemize}

\Cref{tab:confound-grid} applies this checklist to the comparisons available so far. None of AMD's published product comparisons controls every factor. The fixed-frequency prototype comparison controls clock speed but not core count, and the server comparison controls core count but not frequency. The independent comparisons are graded the same way below the line.

\begin{table}[htbp]
  \centering
  \caption{Which confounding factors each available comparison controls (\checkmark{} = controlled, $\times$ = not controlled, $\sim$ = estimated and normalised, ? = not stated; AMD-reported rows above the line, independent below)}
  \label{tab:confound-grid}
  \small
  \begin{tabular}{@{}p{5.6cm}cccc@{}}
    \toprule
    \textbf{Comparison} & \textbf{Clock} & \textbf{Cores} & \textbf{Memory} & \textbf{Same core design} \\
    \midrule
    12-core 3D prototype vs Ryzen~9 5900X, fixed frequency \cite{swaminathanAdvancedPackagingEnabling2021} & \checkmark & \checkmark & ? & \checkmark \\
    Ryzen~7 5800X3D vs Ryzen~9 5900X \cite[slides~29, 35]{wuu2022vcache} & $\times$ & $\times$ & \checkmark & \checkmark \\
    16-core EPYC with V-Cache vs EPYC~73F3 \cite[slides~31, 35]{wuu2022vcache} & $\times$ & \checkmark & ? & \checkmark \\
    Zen to Zen~2 IPC \cite{suggs2019zen2} & ? & ? & ? & $\times$ \\
    Zen~2 to Zen~3 IPC, 3800XT vs 5800X at 4\,GHz \cite{advancedmicrodevicesinc.AMDLaunchesAMD2020} & \checkmark & \checkmark & \checkmark & $\times$ \\
    \midrule
    7950X3D V-Cache CCD vs plain CCD, both at 4.2\,GHz \cite{lamAMDs7950X3DZen2023} & \checkmark & \checkmark & \checkmark & \checkmark \\
    EPYC 7773X vs 7763 \cite{lawendaEvaluatingImpactL32025,domkeLocusPerformanceQuantifying2023} & $\times$ & \checkmark & \checkmark & \checkmark \\
    Milan-X vs Milan, one CCD \cite{cozmaDeepDivingZen2022} & $\sim$ & \checkmark & ? & \checkmark \\
    Ryzen~7 5800X3D vs 5800X \cite{larabelAMDRyzen72022} & $\times$ & \checkmark & ? & \checkmark \\
    Ryzen~7 5800X vs 3700X \cite{waltonAMDRyzen72020} & $\times$ & \checkmark & \checkmark & $\times$ \\
    \bottomrule
  \end{tabular}
\end{table}

The prototype comparison holds both clock (4.0\,GHz) and core count (12) constant but does not state the memory configuration \cite[slide~19]{swaminathanAdvancedPackagingEnabling2021}. Among the independent results, Lam's per-CCD measurements control every listed factor, which makes them the strongest available evidence for the effect of the stacked cache itself \cite{lamAMDs7950X3DZen2023}.

\subsection{Worked example: correcting for clock speed}
\label{sec:clock-correction}
Where clocks were not fixed, a rough correction is possible if the clocks during the test are known. If performance scaled linearly with frequency, dividing a measured speedup by the ratio of clock speeds would give the per-clock effect. Memory-bound code does not scale linearly, because memory latency in nanoseconds does not improve with core clock, so the true per-clock effect lies between the raw ratio and the corrected one. \Cref{tab:clock-correction} applies this to the comparisons whose clocks were measured or stated.

\begin{table}[htbp]
  \centering
  \caption{Worked example: speedups bounded by a clock correction (corrected = measured $\div$ clock ratio)}
  \label{tab:clock-correction}
  \small
  \begin{tabularx}{\textwidth}{@{}>{\raggedright\arraybackslash}X>{\raggedright\arraybackslash}p{3.2cm}cc@{}}
    \toprule
    \textbf{Comparison} & \textbf{Clock ratio (V-Cache / plain)} & \textbf{Measured} & \textbf{Clock-corrected} \\
    \midrule
    OpenFOAM, best mesh \cite{lawendaEvaluatingImpactL32025} & $2.9/3.05 \approx 0.95$ (midpoints of measured ranges) & 1.54$\times$ & 1.62$\times$ \\
    OpenFOAM, worst mesh \cite{lawendaEvaluatingImpactL32025} & $\approx 0.95$ & 0.81$\times$ & 0.85$\times$ \\
    gem5, one CCD \cite{cozmaDeepDivingZen2022} & $3.15/3.3 \approx 0.95$ (estimated) & 1.076$\times$ & 1.13$\times$ \\
    MiniFE, smallest input \cite{domkeLocusPerformanceQuantifying2023} & $2.20/2.45 \approx 0.90$ (base clocks) & 0.89$\times$ & 0.99$\times$ \\
    \bottomrule
  \end{tabularx}
\end{table}
\Cref{fig:clock-correction} plots each comparison before and after the correction.

\input{"../initial-draft-2/figures/clock-correction.tex"}


The correction changes the picture in two ways. For MiniFE's smallest input the whole 11\% loss disappears, which supports the reading that Milan-X loses there only because of its clock. For OpenFOAM's worst mesh a loss of about 15\% remains even after correction, so the clock alone does not explain it; Lawenda \emph{et al.} do not analyse that case, and the four extra cycles of L3 latency cannot account for so large a difference on their own. This is a question the published evidence leaves open. Because base clocks and measured ranges are crude stand-ins for the actual frequency during each run, the corrected values are indicative only.

\section{Summary of Part V}
The evidence answers RQ4 as follows. The stacked L3 costs about four cycles of L3 latency in every generation measured, and in its first two generations also several percent of clock speed. In exchange it removes capacity misses for working sets between 32 and 96\,MB per CCD. Where those misses dominate, the benefit is large: from about 10--20\% higher IPC in miss-heavy games to more than 1.5$\times$ in memory-bound simulation, and up to 3.4$\times$ for one HPC proxy at favourable problem sizes. Where they do not, the benefit is small, absent or slightly negative. Separating the cache effect from other factors requires either fixed-clock comparisons or performance-counter evidence; most vendor comparisons provide neither, and the best-controlled independent study available covers only a handful of workloads. The confidence of these conclusions is therefore high for the direction and mechanism, and moderate for the size of the effect in any particular workload.
